\documentclass[12pt]{article}
\usepackage[utf8]{inputenc}
\usepackage[T1]{fontenc}
\usepackage[margin=1in]{geometry}
\usepackage{setspace}
\usepackage{xcolor}
\usepackage[colorlinks=true, urlcolor=blue!60!black, linkcolor=blue!60!black]{hyperref}
\setlength{\parindent}{0em}
\setlength{\parskip}{0.45em}
\renewcommand{\baselinestretch}{1.0}
\begin{document}
\begin{center}
{\Large \textit{United States v. Comstock}, 560 U.S. 126 (2010)\par}
\vspace{0.3em}
{\normalsize Case Brief --- One-Page Edition\par}
\end{center}

\textbf{Memory Jogger.} The sexually-dangerous-prisoner case: civil commitment as a necessary and proper incident of federal custody (\href{https://supreme.justia.com/cases/federal/us/560/126/}{560 U.S. 126}).

\textbf{Facts.} 18 U.S.C. \S~4248 lets federal courts commit sexually dangerous persons in federal custody, including prisoners nearing release, unless a State takes custody.

\textbf{Procedural History.} The District Court and the Fourth Circuit held that Congress lacked Article I power.

\textbf{Votes.} Reversed, 7--2. Breyer wrote for the Court; Kennedy and Alito concurred in the judgment; Thomas and Scalia dissented.

\textbf{Holding.} The Necessary and Proper Clause authorizes \S~4248 as a means rationally related to implementing enumerated powers.

\textbf{Reasoning.} The Clause is broad under \textit{McCulloch}; Congress has long provided for mentally ill federal prisoners; a responsible custodian may prevent dangerous releases; the statute defers to States; and it reaches only people already in federal custody. Kennedy and Alito demanded a stronger link than ``rational basis''; Thomas found the power behind a conviction exhausted at sentence's end.

\textbf{Examples.} \emph{Same-Side:} isolating a contagious inmate past release. \emph{Opposite-Side:} committing state felons never in federal custody. \emph{Fence-Sitter:} registration for long-released offenders.

\textbf{Critique.} Stacking implied powers leaves no clear stopping point.

\textbf{Key Quotation.} ``Taken together, these considerations lead us to conclude that the statute is a `necessary and proper' means of exercising the federal authority that permits Congress to create federal criminal laws, to punish their violation, to imprison violators, to provide appropriately for those imprisoned, and to maintain the security of those who are not imprisoned but who may be affected by the federal imprisonment of others.'' --- Justice Breyer, Opinion of the Court, 560 U.S. at 149.

\end{document}
